\chapter{Referencial Teórico} \label{Cap:ReferencialTeorico}

\FloatBarrier
\section{Customer Churn: Conceituação e Importância Estratégica} \label{sec:churn-conceito}

O conceito de customer churn refere-se ao processo pelo qual clientes encerram ou interrompem seu relacionamento comercial com uma organização, deixando de realizar compras ou de utilizar seus serviços \cite{manzoor2024}. No contexto do comércio eletrônico, o churn pode manifestar-se de forma explícita, como no cancelamento de assinaturas, ou implícita, caracterizada pela ausência de novas transações em um determinado período, o que torna sua identificação mais complexa e dependente da definição de janelas temporais de observação.

Do ponto de vista estratégico, a gestão do churn ocupa posição central nos modelos de Customer Relationship Management (CRM), uma vez que a rentabilidade de longo prazo das organizações está diretamente vinculada à capacidade de manutenção e expansão de sua base de clientes. Estudos clássicos demonstraram que reduções de 5\% na taxa de churn podem resultar em incrementos de 25\% a 95\% nos lucros, em função do aumento do valor gerado pelo cliente ao longo do ciclo de vida do relacionamento \cite{reichheld2000}. Adicionalmente, a geração de receita por clientes recorrentes tende a ser mais estável e previsível do que aquela proveniente de novos clientes, cujo custo de aquisição é substancialmente superior ao de retenção, especialmente no comércio eletrônico \cite{reichheld2000}.

No setor de e-commerce, o desafio da retenção é amplificado pela abundância de concorrentes e pela facilidade com que os consumidores podem migrar para outras plataformas. Como as taxas de churn variam consideravelmente entre setores \cite{imani2025}, torna-se imperativa a adoção de estratégias preditivas que permitam antecipar o abandono e direcionar ações de retenção de forma proativa.

\FloatBarrier
\section{Metodologia RFM e Engenharia de Features} \label{sec:rfm}

A metodologia RFM, acrônimo para Recência (R), Frequência (F) e Monetário (M), constitui um dos frameworks mais consolidados para segmentação e análise comportamental de clientes, com origem na área de marketing direto \cite{hughes1994}. A Recência mensura o intervalo de tempo desde a última transação do cliente; a Frequência quantifica o número de transações realizadas em um período determinado; e o valor Monetário corresponde ao volume financeiro total gasto pelo cliente no período de análise.

No contexto de modelos preditivos de churn, o framework RFM serve como base para a engenharia de features, uma vez que as três dimensões capturam de forma objetiva os principais aspectos do comportamento transacional do cliente. Clientes com elevada recência, ou seja, com longa data desde a última compra, tendem a apresentar maior propensão ao churn, enquanto aqueles com alta frequência e alto valor monetário representam os perfis mais valiosos e tipicamente mais estáveis.

Um aspecto metodológico crítico na construção de features RFM para modelos preditivos é o tratamento rigoroso da dimensão temporal, de modo a evitar o vazamento de informação futura para o conjunto de treino, problema conhecido como data leakage \cite{kaufman2012}. Uma estratégia consolidada para mitigar esse risco é a separação entre uma janela de observação, na qual as features são computadas, e uma janela de predição, na qual o target é definido, estratégia detalhada na Seção \ref{sec:janela-temporal}.

\FloatBarrier
\section{Algoritmos de Machine Learning para Classificação} \label{sec:algoritmos}

A literatura sobre predição de churn contempla ampla variedade de algoritmos de classificação, com desempenhos variáveis a depender das características do dataset, do grau de desbalanceamento entre classes \cite{he2009} e da complexidade das relações não lineares entre as features \cite{vafeiadis2015}. Entre os modelos mais amplamente aplicados, destacam-se a Regressão Logística, os métodos baseados em árvores e os algoritmos ensemble.

A Regressão Logística, modelo linear de referência amplamente utilizado em estudos de churn por sua interpretabilidade e eficiência computacional, apresenta limitações diante de relações não lineares entre variáveis, sendo frequentemente superada por modelos baseados em árvores em contextos de maior complexidade \cite{prabadevi2023}.

O Random Forest consiste em um ensemble de árvores de decisão treinadas com bootstrap sampling e seleção aleatória de features, o que confere ao modelo robustez ao overfitting e capacidade de capturar interações não lineares entre variáveis \cite{breiman2001}. O algoritmo fornece, naturalmente, medidas de importância de features via impureza de Gini, tornando-se especialmente útil em contextos onde a interpretabilidade do modelo é relevante.

Os métodos de Gradient Boosting, incluindo o Gradient Boosting clássico \cite{friedman2001} e suas variantes otimizadas como o HistGradient Boosting, operam pela construção sequencial de modelos fracos (weak learners), cada um focado na correção dos erros do anterior, minimizando uma função de perda de forma iterativa. Tais abordagens têm consistentemente demonstrado desempenho superior em benchmarks de classificação com dados tabulares.

\FloatBarrier
\section{Desbalanceamento de Classes e Técnicas de Balanceamento} \label{sec:desbalanceamento}

O desbalanceamento de classes constitui um dos principais desafios práticos em problemas de predição de churn, especialmente no contexto do e-commerce, onde a proporção de clientes que efetivamente retornam é tipicamente muito menor do que a daqueles que não voltam a comprar. Quando não tratado adequadamente, o desbalanceamento induz os modelos a favorecerem a classe majoritária, resultando em alta acurácia geral, porém com baixa capacidade de identificação dos casos de churn \cite{he2009}.

A literatura propõe diversas abordagens para mitigar esse problema. O oversampling da classe minoritária, como implementado pelo algoritmo SMOTE (Synthetic Minority Over-sampling Technique), gera instâncias sintéticas baseadas na interpolação entre exemplos reais da classe minoritária \cite{chawla2002}. O undersampling da classe majoritária, por sua vez, reduz a proporção dos exemplos mais frequentes, equilibrando o dataset de treino. Estudos indicam que a combinação de ambas as estratégias tende a produzir resultados superiores às abordagens isoladas: técnicas híbridas como SMOTE combinado com Tomek Links ou com Edited Nearest Neighbors (ENN) geram conjuntos de dados com fronteiras de classe mais bem definidas, atenuando o overfitting \cite{batista2004}. No domínio específico da predição de churn, estudos comparativos recentes avaliam o desempenho de diferentes métodos de resampling sobre dados altamente desbalanceados provenientes de setores como telecomunicações, varejo online e bancário \cite{haddadi2024}.

\FloatBarrier
\section{Métricas de Avaliação para Dados Desbalanceados} \label{sec:metricas}

A seleção de métricas de avaliação adequadas é uma decisão crítica em problemas de classificação com dados desbalanceados. A acurácia global constitui métrica enganosa nesses contextos, uma vez que um classificador trivial que prediz sempre a classe majoritária pode atingir alta acurácia sem qualquer capacidade preditiva real. De la Cruz Huayanay, Bazán e Russo (\citeyear{delacruz2025}) demonstraram que a ROC-AUC isolada é insuficiente como critério de seleção de modelos em dados desbalanceados, ao mostrar, em estudo de simulação com modelo verdadeiro conhecido, que ela não distingue o modelo corretamente especificado do mal especificado em nenhum cenário avaliado, enquanto MCC, G-Mean e Kappa de Cohen o fazem de forma consistente: achado que motivou a adoção de um portfólio de métricas complementares neste trabalho.

Entre as métricas usualmente empregadas para contextos desbalanceados, destacam-se cinco. A ROC-AUC (Area Under the ROC Curve) mede a capacidade discriminatória geral do modelo, de forma independente do threshold \cite{fawcett2006}. O G-Mean (Geometric Mean), proposto originalmente por \citeonline{kubat1997}, é definido como a raiz quadrada do produto entre Sensibilidade (Recall da classe positiva) e Especificidade (Recall da classe negativa), penalizando modelos que sacrificam o desempenho em uma das classes. O MCC (Coeficiente de Correlação de Matthews), introduzido por \citeonline{matthews1975}, é considerado por \citeonline{chicco2020} a métrica mais informativa e confiável para classificação binária, por levar em conta, de forma simétrica, todos os quadrantes da matriz de confusão (verdadeiros positivos, verdadeiros negativos, falsos positivos e falsos negativos). O Kappa de Cohen \cite{cohen1960} corrige o desempenho pelo acerto esperado ao acaso. Por fim, o F1-Macro, média harmônica não ponderada de Precisão e Recall por classe, complementa o conjunto ao capturar o equilíbrio entre as classes no ponto de corte adotado. \citeonline{luque2019} demonstram analiticamente como o desbalanceamento de classes distorce o comportamento dessas métricas, o que reforça a necessidade de interpretá-las em conjunto e à luz da distribuição real do conjunto de teste.

A Curva Precision-Recall (PR), especialmente sua área (AUPRC ou Average Precision), é preferível à ROC-AUC em cenários de desbalanceamento severo, por ser mais sensível ao desempenho na classe minoritária e não ser afetada pela grande quantidade de verdadeiros negativos típica desses contextos \cite{saito2015}. Cabe registrar que a Average Precision pondera de forma desproporcional o extremo inicial da ordenação, região em que a precisão é estimada sobre poucas observações, característica que se torna relevante na interpretação conduzida na Seção \ref{sec:benchmark}.

\FloatBarrier
\section{Explicabilidade de Modelos: Valores SHAP e Importância de Features} \label{sec:shap-teorico}

A crescente adoção de modelos complexos de Machine Learning em contextos de negócio impõe desafios à interpretabilidade das predições, especialmente em organizações onde decisões baseadas em algoritmos precisam ser justificadas a gestores e reguladores. Nesse contexto, a área de Inteligência Artificial Explicável (XAI) tem ganhado relevância crescente, propondo métodos para tornar os modelos mais transparentes sem necessariamente sacrificar seu desempenho.

Os valores SHAP (SHapley Additive exPlanations) foram propostos por \citeonline{lundberg2017} como um framework unificado de explicabilidade, fundamentado na teoria dos jogos cooperativos de Shapley \cite{shapley1953} e dotado de propriedades teóricas desejáveis, como consistência e atribuição aditiva. Desde então, tornaram-se um dos métodos mais consolidados para explicabilidade local e global de modelos de Machine Learning. Para cada predição, os valores SHAP quantificam a contribuição marginal de cada feature para o desvio do resultado em relação ao valor base do modelo. O TreeSHAP, implementado na classe TreeExplainer da biblioteca shap \cite{lundberg2020}, permite o cálculo eficiente de valores SHAP para modelos baseados em árvores, como o Random Forest. Cabe registrar que o TreeSHAP decompõe as predições percorrendo a mesma estrutura de árvores ajustada da qual deriva a importância por impureza de Gini, de modo que a concordância entre os dois estimadores é esperada por construção e não constitui, isoladamente, validação independente da hierarquia de importância obtida.

Para modelos lineares, a mesma decomposição admite forma fechada. Sob a hipótese de independência entre as features, o valor SHAP da variável $j$ na instância $x$ é dado por $\varphi_j(x) = \beta_j \cdot (x_j - \mathbb{E}[x_j])$, em que $\beta_j$ é o coeficiente estimado. Essa propriedade permite que os valores SHAP sejam calculados para estimadores de naturezas distintas, o que os torna o único terreno comum de comparação de importância entre modelos que não compartilham medida nativa, recurso empregado na Seção \ref{sec:interpretabilidade}. Cabe a ressalva de que as contribuições assim obtidas encontram-se em escalas diferentes conforme a saída do modelo: em espaço de probabilidade para o Random Forest sob TreeSHAP, e em log-odds para modelos lineares e para o boosting, de modo que apenas posições e participações relativas são comparáveis entre eles.

\FloatBarrier
\section{Estudos Relacionados} \label{sec:estudos-relacionados}

A predição de churn em e-commerce tem sido objeto de crescente atenção na literatura científica. Revisões clássicas do estado da arte já apontavam o bom desempenho de árvores de decisão e redes neurais em relação a modelos de regressão tradicionais para predição de churn \cite{hadden2007}.

No campo específico da avaliação de modelos sob desbalanceamento, De la Cruz Huayanay, Bazán e Russo (\citeyear{delacruz2025}) conduziram estudo de simulação com doze métricas de classificação, concluindo que o MCC, o G-Mean e o Kappa de Cohen mantêm desempenho consistente na identificação do modelo corretamente especificado, enquanto a AUC e a acurácia falham em todos os cenários avaliados. Os autores registram duas delimitações de alcance que dialogam diretamente com este trabalho: o estudo não contempla amostras pequenas, tendo empregado tamanhos de 5.000 e 10.000 observações, e sugerem, como extensão, a comparação com algoritmos de classificação como Naive Bayes, K-Nearest Neighbors, Árvore de Decisão, Support Vector Machines e Random Forest. O presente trabalho situa-se nesse recorte, ao avaliar esses algoritmos sobre uma base de 1.178 clientes com aproximadamente cinquenta observações na classe minoritária.

Revisões sistemáticas recentes mapearam tendências e desafios da área. \citeonline{imani2025}, ao sintetizarem estudos publicados entre 2020 e 2024, documentam a ampla adoção de métodos ensemble baseados em árvores, que mantêm desempenho competitivo em datasets tabulares de tamanho moderado, contexto no qual modelos mais complexos nem sempre compensam o custo computacional adicional. \citeonline{manzoor2024} complementam esse panorama ao mapear 212 estudos publicados entre 2015 e 2023, sistematizando recomendações de engenharia de features e de avaliação para modelos de churn. Os autores destacam a importância de variáveis comportamentais complementares ao RFM e enfatizam que o desempenho preditivo depende fortemente da qualidade da preparação dos dados, e não apenas da escolha do algoritmo. Em um recorte mais específico, \citeonline{haddadi2024} conduziram estudo comparativo sobre técnicas de resampling para churn em datasets desbalanceados, avaliando métodos como SMOTE, ADASYN e uma abordagem de resampling em duas fases, esta última especialmente compatível com redes neurais profundas.

No contexto específico do dataset Olist \cite{olist2018}, \citeonline{matuszelanski2022} conduziram um dos estudos mais completos disponíveis, combinando abordagens espaciais e de Machine Learning. Os autores documentaram a característica predominante de clientes com compra única como desafio metodológico central e recomendaram filtros baseados na frequência mínima de compras. A adoção desses filtros, aliada à definição rigorosa de janelas temporais, constitui estratégia necessária para tornar o problema tratável do ponto de vista de Machine Learning, conforme proposto neste trabalho.

\FloatBarrier
